\documentclass[10pt,a4paper]{amsart}
\usepackage{amsmath,amssymb,amsthm,graphicx}
\usepackage{verbatim}
\newenvironment{reference}{\comment}{\endcomment}
%\newenvironment{reference}{}{}
\newenvironment{slogan}{\comment}{\endcomment}
\newenvironment{history}{\comment}{\endcomment}
% For commutative diagrams we use Xy-pic
\usepackage[all]{xy}

% We use 2cell for 2-commutative diagrams.
\xyoption{2cell}
% Theorem environments.
%
\theoremstyle{plain}
\newtheorem{theorem}[subsection]{Theorem}
\newtheorem{proposition}[subsection]{Proposition}
\newtheorem{lemma}[subsection]{Lemma}

\theoremstyle{definition}
\newtheorem{definition}[subsection]{Definition}
\newtheorem{example}[subsection]{Example}
\newtheorem{exercise}[subsection]{Exercise}
\newtheorem{situation}[subsection]{Situation}

\theoremstyle{remark}
\newtheorem{remark}[subsection]{Remark}
\newtheorem{remarks}[subsection]{Remarks}

\numberwithin{equation}{subsection}

% Macros
%
\def\lim{\mathop{\mathrm{lim}}\nolimits}
\def\colim{\mathop{\mathrm{colim}}\nolimits}
\def\Spec{\mathop{\mathrm{Spec}}}
\def\Hom{\mathop{\mathrm{Hom}}\nolimits}
\def\Ext{\mathop{\mathrm{Ext}}\nolimits}
\def\SheafHom{\mathop{\mathcal{H}\!\mathit{om}}\nolimits}
\def\SheafExt{\mathop{\mathcal{E}\!\mathit{xt}}\nolimits}
\def\Sch{\mathit{Sch}}
\def\Mor{\mathop{\mathrm{Mor}}\nolimits}
\def\Ob{\mathop{\mathrm{Ob}}\nolimits}
\def\Sh{\mathop{\mathit{Sh}}\nolimits}
\def\NL{\mathop{N\!L}\nolimits}
\def\CH{\mathop{\mathrm{CH}}\nolimits}
\def\proetale{{pro\text{-}\acute{e}tale}}
\def\etale{{\acute{e}tale}}
\def\QCoh{\mathit{QCoh}}
\def\Ker{\mathop{\mathrm{Ker}}}
\def\Im{\mathop{\mathrm{Im}}}
\def\Coker{\mathop{\mathrm{Coker}}}
\def\Coim{\mathop{\mathrm{Coim}}}

% Boxtimes
%
\DeclareMathSymbol{\boxtimes}{\mathbin}{AMSa}{"02}

%
% Macros for moduli stacks/spaces
%
\def\QCohstack{\mathcal{QC}\!\mathit{oh}}
\def\Cohstack{\mathcal{C}\!\mathit{oh}}
\def\Spacesstack{\mathcal{S}\!\mathit{paces}}
\def\Quotfunctor{\mathrm{Quot}}
\def\Hilbfunctor{\mathrm{Hilb}}
\def\Curvesstack{\mathcal{C}\!\mathit{urves}}
\def\Polarizedstack{\mathcal{P}\!\mathit{olarized}}
\def\Complexesstack{\mathcal{C}\!\mathit{omplexes}}
% \Pic is the operator that assigns to X its picard group, usage \Pic(X)
% \Picardstack_{X/B} denotes the Picard stack of X over B
% \Picardfunctor_{X/B} denotes the Picard functor of X over B
\def\Pic{\mathop{\mathrm{Pic}}\nolimits}
\def\Picardstack{\mathcal{P}\!\mathit{ic}}
\def\Picardfunctor{\mathrm{Pic}}
\def\Deformationcategory{\mathcal{D}\!\mathit{ef}}
\usepackage{cleveref}
\crefname{originaldefinition}{original definition}{original definitions}
\Crefname{originaldefinition}{Original definition}{Original definitions}
\makeatletter
\newlabel{definition-catenary}{{definition-catenary}{source}{}{sourcepointer-1}{}}
\newlabel{definition-catenary@cref}{{[originaldefinition][0][]definition-catenary}{[1][0][]source}{}{}{}}
\makeatother
\title{Catenarity is detected on chains of a fixed codimension}
\author{The Stacks Project Authors}
\date{stacks-project:topology.tex:local-snapshot-20261012; GNU FDL 1.2 or later}
\begin{document}
\maketitle
\noindent\textit{Editorial note.} All delivered mathematical text is the author's own, copied from the pinned Stacks Project chapter source sources/stacks-topology/pinned-original-source.tex; no statement or proof step was rewritten, reordered or removed.\par\medskip
\noindent\textit{Editorial note.} Preamble provenance: the Stacks Project keeps one single shared preamble for every chapter; this card takes it directly from the staged sources/stacks-topology/preamble.tex as three verbatim spans (lines 8-12, 14-18 and 154-217), so no derived file was written and nothing was added or redefined. The wrapper class and the base packages are supplied by the assembler.\par\medskip
\noindent\textit{Editorial note.} Bibliography provenance: neither the counted statement nor its proof contains a citation (this lane rescans the whole span before writing the card), so the card carries no bibliography and no entry was reconstructed.\par\medskip
\noindent\textit{Editorial note.} Selection provenance: this card was opened under the file-based control contract (CONTROL.md), with the mechanical pre-card gate run on its statement and proof interval before the card was written.\par\medskip
\noindent\textit{Allowed standard prerequisites (editorial).}\begin{itemize}
\item The definition of a catenary space through equality of dimensions of maximal chains. Used inside the counted proof (source lines 1797-1810). Reduces catenarity to the equality of the dimensions of two chains.
\end{itemize}\medskip
\ifdefined\hypertarget\hypertarget{sourcepointer-1}{}\fi
\noindent\textit{Source reference.} Original definition definition-catenary: sources/stacks-topology/pinned-original-source.tex, source lines 1752--1752. Navigation pointer only: the counted text refers to the Stacks item definition-catenary, which is an existing labeled result of this same chapter and is not part of this card. The pointer carries the target's own label as its number and its definition line as its span; its content is not reproduced here.\par\medskip
\begin{lemma}
\label{lemma-catenary-in-codimension}
Let $X$ be a topological space. The following are equivalent:
\begin{enumerate}
\item $X$ is catenary, and
\item for every pair of irreducible closed subsets $Y \subset Y'$ we have
$\text{codim}(Y, Y') < \infty$ and for every triple
$Y \subset Y' \subset Y''$ of irreducible closed subsets we have
$$
\text{codim}(Y, Y'') = \text{codim}(Y, Y') + \text{codim}(Y', Y'').
$$
\end{enumerate}
\end{lemma}

\begin{proof}
Let suppose that $X$ is catenary. According to
Definition \ref{definition-catenary}, for every pair of
irreducible closed subsets $Y \subset Y'$ we have
$\text{codim}(Y,Y') < \infty$. Let $Y \subset Y' \subset Y''$
be a triple of irreducible closed subsets of $X$. Let
$$
Y = Y_0 \subset Y_1 \subset ... \subset Y_{e_1} = Y'
$$
be a maximal chain of irreducible closed subsets between $Y$ and $Y'$
where $e_1 = \text{codim}(Y,Y')$. Let also
$$
Y' = Y_{e_1} \subset Y_{e_1 + 1}\subset ... \subset Y_{e_1 + e_2} = Y''
$$
be a maximal chain of irreducible closed subsets between $Y'$
and $Y''$ where $e_2 = \text{codim}(Y',Y'')$.
As the two chains are maximal, the concatenation
$$
Y = Y_0\subset Y_1 \subset ... \subset Y_{e_1} = Y' =
Y_{e_1} \subset Y_{e_1+1}\subset ... \subset Y_{e_1+e_2}=Y''
$$
is maximal too (between $Y$ and $Y''$) and its length equals to $e_1 + e_2$.
As $X$ is catenary, each maximal chain has the same length equals to the
codimension. Thus the point (2) that
$\text{codim}(Y,Y'') = e_1 + e_2 = \text{codim}(Y,Y') + \text{codim}(Y',Y'')$
is verified.

\medskip\noindent
For the reciprocal, we show by induction that : if $Y = Y_1 \subset ... \subset
Y_n = Y'$, then $ \text{codim}(Y,Y') = \text{codim}(Y_1,Y_2) + ... +
\text{codim}(Y_{n-1},Y_n)$. Therefore, it forces maximal chains to have the
same length.
\end{proof}

\end{document}
